% =============================================================================
% Presentación para reunión con Gonzalo — resultados y plan de publicación
% Compilar: pdflatex presentacion_gonzalo.tex  (dos pasadas)
% =============================================================================
\documentclass[11pt,aspectratio=169]{beamer}

\usetheme{Madrid}
\usecolortheme{seahorse}
\setbeamertemplate{navigation symbols}{}
\setbeamertemplate{footline}[frame number]

\usepackage[utf8]{inputenc}
\usepackage[T1]{fontenc}
\usepackage[spanish,es-nodecimaldot]{babel}
\usepackage{amsmath,amssymb}
\usepackage{graphicx}
\usepackage{booktabs}

\newcommand{\eps}{\varepsilon}
\newcommand{\gstar}{g^*}
\newcommand{\lammps}{\textsc{Lammps}}

\title[Rebote de nanocúmulos LJ]{¿A dónde se va la energía de una pelota
que rebota?\\ \large Partición energética del impacto de nanocúmulos LJ
--- resultados y plan de publicación}
\author[F. González-Wasaff]{Felipe González-Wasaff\\
\footnotesize reunión con Gonzalo Gutiérrez}
\institute[IMA-PUCV]{Instituto de Matemáticas, PUCV}
\date{Julio 2026}

\begin{document}

% -----------------------------------------------------------------------------
\begin{frame}
\titlepage
\end{frame}

% -----------------------------------------------------------------------------
\begin{frame}{De dónde venimos y qué hicimos}
\begin{columns}[T]
\column{0.55\textwidth}
\textbf{Punto de partida (2012):} capítulo 4 del libro
\emph{Inelastic Collisions and Hypervelocity Impacts at Nanoscopic Level}
--- el rebote de la esfera de $N\approx79$ átomos de Ar.

\bigskip
\textbf{Lo que hicimos ahora:}
\begin{itemize}
    \item Reproducción completa con \lammps{} (código abierto, repo público)
    \item Cúmulos pre-minimizados de $N$ \emph{exacto}
    \item \textbf{Partición al primer apogeo en 4 canales} con cierre exacto
    \item Barrido $N\in\{1,2,7,13,55,79,87\}$ + equipartición
    \item Simulación larga ($5000\,\tau$) + caso con adhesión
\end{itemize}
\column{0.42\textwidth}
\begin{block}{La pregunta del paper}
El coeficiente de restitución $e<1$ dice \emph{cuánta} energía no vuelve.
\medskip

¿\textbf{Dónde} queda esa energía justo después del rebote?
¿En la pelota o en el piso?
\end{block}
\end{columns}
\end{frame}

% -----------------------------------------------------------------------------
\begin{frame}{El resultado en una imagen}
\begin{center}
    \includegraphics[height=0.82\textheight]{../data/outputs/fig_sequence.png}
\end{center}
\end{frame}

% -----------------------------------------------------------------------------
\begin{frame}{Los cuatro canales ($N=79$, repulsivo, $T=0$)}
\begin{columns}[T]
\column{0.58\textwidth}
\small
\begin{tabular}{cccccc}
\toprule
$h_0^{\rm ef}$ & $v_{\rm imp}$ & $\alpha_{\rm cm}$ & $\alpha_{T_B}$ &
$\alpha_{\Delta V_B}$ & $\alpha_{\rm surf}$ \\
\midrule
 8.8 & 0.68 & 0.853 & 0.021 & 0.056 & 0.070 \\
17.2 & 0.95 & 0.502 & 0.213 & 0.208 & 0.077 \\
25.9 & 1.16 & 0.304 & 0.319 & 0.310 & 0.066 \\
41.8 & 1.47 & 0.401 & 0.257 & 0.264 & 0.078 \\
56.9 & 1.72 & 0.378 & 0.197 & 0.328 & 0.098 \\
79.0 & 2.03 & 0.345 & 0.204 & 0.294 & 0.156 \\
\bottomrule
\end{tabular}

\medskip
\footnotesize $\alpha_{\rm cm}+\alpha_{T_B}+\alpha_{\Delta V_B}
+\alpha_{\rm surf}=1$; cierre medido $<10^{-4}$.
\column{0.40\textwidth}
\begin{itemize}
    \item Plateau $\alpha_{T_B}\approx0.22$ para $v>0.9\,\sigma/\tau$
    \item Canal nuevo: \textbf{expansión anarmónica}
    $\alpha_{\Delta V_B}\approx0.2$--$0.3$ (comparable al cinético)
    \item La superficie se lleva solo 7--16\%
\end{itemize}
\begin{alertblock}{\small Mensaje central}
\small La energía ``perdida'' queda mayoritariamente
\textbf{dentro de la bola} ($\alpha_{\rm bola}\approx0.4$--$0.6$),
no en el piso.
\end{alertblock}
\end{columns}
\end{frame}

% -----------------------------------------------------------------------------
\begin{frame}{Validación contra 2012 (tu Fig.~7, digitalizada)}
\begin{columns}[T]
\column{0.52\textwidth}
\small
\begin{tabular}{c|ccc|cc}
\toprule
 & \multicolumn{3}{c|}{2012 (Fig.~7)} & \multicolumn{2}{c}{2026} \\
$h_0$ & $h_0^{\rm ef}$ & $h_1^{\rm ef}$ & $e_1$ & $h_0^{\rm ef}$ & $e_1$ \\
\midrule
11.03 &  0.6 &  0.4 & ---  &  8.8 & 0.92 \\
18.64 &  8.4 &  6.8 & 0.90 & 17.2 & 0.71 \\
28.15 & 17.9 &  3.3 & 0.43 & 25.9 & 0.55 \\
41.47 & 31.3 &  8.2 & 0.51 & 41.8 & 0.63 \\
58.60 & 48.4 & 14.3 & 0.54 & 56.9 & 0.61 \\
79.53 & 69.2 & 18.2 & 0.51 & 79.0 & 0.59 \\
\bottomrule
\end{tabular}
\column{0.46\textwidth}
\begin{itemize}
    \item \textbf{Hallazgo de convención:} en 2012, $h_0$ = altura
    \emph{absoluta} del CM ($z_{\rm rest}\approx10.3\,\sigma$);
    ahora se mide desde el tope de la caja
    \item A caída efectiva comparable: \textbf{0.90 vs 0.92} en el
    régimen cuasi-elástico
    \item A alta energía, 2012 queda $\sim$0.06--0.09 abajo
    (cúmulo sin minimizar, $N$ variable 79--84, Beeman)
\end{itemize}
\end{columns}
\end{frame}

% -----------------------------------------------------------------------------
\begin{frame}{Un punto de autoría que necesito confirmar contigo}
\begin{block}{Lo que el capítulo 2012 establece}
Canales $K_B, T_B, V_B, U_c, \ldots$; alturas por rebote (Fig.~7);
$U_c$ domina el balance \emph{durante} el contacto; fusión de la bola
vía $g(r)$ para $h_0>41.5\,\sigma$.
\end{block}
\begin{alertblock}{Lo que el borrador presenta como contribución NUEVA (2026)}
\begin{itemize}
    \item El plateau cuantificado $\alpha_{T_B}\approx0.22$ al apogeo
    \item La fórmula de equipartición
    $\alpha_{\rm int}=(3N-3)/(3N-3+3N_c)$
    \item La escala $h^*=T_B/(Mg)$
\end{itemize}
\end{alertblock}
\centering
¿De acuerdo con esa atribución? Si algo viene de notas tuyas no
publicadas, lo citamos como corresponde.
\end{frame}

% -----------------------------------------------------------------------------
\begin{frame}{Equipartición: el tamaño de la pelota importa}
\begin{center}
    \includegraphics[width=0.9\textwidth]{../data/outputs/fig_alpha_TB_full_sweep.pdf}
\end{center}
\small Con $\alpha_{\rm bola}$ y $N_c\approx N$:
$\alpha_{\rm bola}\to\tfrac12$ para $N\to\infty$.
Salvedades honestas: datos de $N=79,87$ dan $0.543$ (leve exceso sobre
$1/2$; ¿rotación?); $N=13$ no termaliza (régimen molecular).
\end{frame}

% -----------------------------------------------------------------------------
\begin{frame}{$h^*=T_B/(Mg)$: útil para el primer rebote, no es atractor}
\begin{columns}[T]
\column{0.55\textwidth}
\begin{center}
    \includegraphics[width=\textwidth]{../data/outputs/fig_paso_B_convergencia.pdf}
\end{center}
\column{0.42\textwidth}
\begin{itemize}
    \item Primer rebote: $h_1/h^* = 0.96$ \checkmark
    \item Rebotes sucesivos: la razón cae a $\sim$0.18
    \item Atrapamiento tras 23 rebotes ($t\approx660\,\tau$)
    \item El ``colapso'' posterior a través de la red es un artefacto
    del potencial repulsivo puro (documentado como tal)
\end{itemize}
\end{columns}
\end{frame}

% -----------------------------------------------------------------------------
\begin{frame}{Con adhesión ($c_{ps}=1$): la bola siempre pega}
\begin{columns}[T]
\column{0.55\textwidth}
\begin{center}
    \includegraphics[width=\textwidth]{../data/outputs/fig_paso_C_adhesion.pdf}
\end{center}
\column{0.42\textwidth}
\begin{itemize}
    \item 24/24 casos adheridos ($v$ hasta $2\,\sigma/\tau$)
    \item $E_{\rm adh}\propto N^{0.73}$ ($\sim$ área de contacto)
    \item Umbral extrapolado $h_0^*\sim10^2$--$10^3\,\sigma$
    (solo orden de magnitud: a esas $v$ la bola se funde)
    \item Los grandes escapan más fácil: $h_0^*\propto N^{-0.27}$
    $\to$ conexión con \emph{soft landing}
\end{itemize}
\end{columns}
\end{frame}

% -----------------------------------------------------------------------------
\begin{frame}{Rigor técnico y límites declarados}
\begin{columns}[T]
\column{0.5\textwidth}
\textbf{Resuelto en esta etapa}
\begin{itemize}
    \item \small Potencial truncado \emph{y desplazado}: sin shift, el
    balance fallaba 2--6\% de $E_0$ por impacto; con shift,
    $|\Delta E|<0.03\,\eps$ y cierre $<10^{-4}$
    \item \small $\Delta t=5\times10^{-3}\,\tau$ validado contra la
    conservación
    \item \small Figuras 100\% reproducibles por script; repo público
    con inputs y análisis
\end{itemize}
\column{0.48\textwidth}
\textbf{Límites honestos (= trabajo futuro)}
\begin{itemize}
    \item \small $T=0$, una orientación del cúmulo, una trayectoria por
    punto $\to$ sin barras de error estadísticas todavía
    \item \small No monotonía de $e_1$ con $h_0$: ¿física o varianza?
    \item \small Momento angular post-rebote sin medir
    (\texttt{compute angmom})
    \item \small $N=19,43$ pendientes (transición molécula--sólido)
\end{itemize}
\end{columns}
\end{frame}

% -----------------------------------------------------------------------------
\begin{frame}{Plan editorial: RBEF, sección ``Artigos Gerais''}
\begin{columns}[T]
\column{0.52\textwidth}
\textbf{Manuscrito listo} (\texttt{borrador\_rbef.tex}):
\begin{itemize}
    \item \small Marco didáctico: la pregunta del curso de mecánica,
    métodos como tutorial, 5 actividades propuestas
    \item \small Protocolo RBEF cumplido: bilingüe ES+EN, \texttt{bibitem}
    en orden de citación, sin TikZ, ``ecuación ($n$)''
    \item \small Sin costos de publicación; SciELO, acceso abierto
\end{itemize}
\medskip
\textbf{También existen} (para un eventual segundo paper de
investigación): versión extendida ES + doble columna revtex ES y EN.
\column{0.44\textwidth}
\begin{block}{\small Flujo propuesto}
\small
\begin{enumerate}
    \item Tu revisión del manuscrito
    \item arXiv (\texttt{physics.ed-ph}), simultáneo con
    \item Submissão RBEF (sistema pide ORCID)
    \item Repo GitHub público al someter
\end{enumerate}
\end{block}
\end{columns}
\end{frame}

% -----------------------------------------------------------------------------
\begin{frame}{Decisiones que necesito de ti}
\begin{enumerate}
    \item \textbf{Atribución} de $\alpha_{T_B}\approx0.22$, la fórmula de
    equipartición y $h^*$ (¿contribución nueva o notas previas tuyas?)
    \item \textbf{Idioma} del envío a RBEF: ¿español (aceptado) o
    portugués/inglés?
    \item Tu \textbf{ORCID} (y visto bueno del flujo arXiv $\to$ RBEF)
    \item \textbf{Autor correspondiente} y orden de autores
    \item ¿Guardamos la versión de investigación (doble columna) para una
    revista de física computacional más adelante, con la estadística de
    orientaciones incluida?
\end{enumerate}
\bigskip
\centering
\footnotesize Repo: \texttt{github.com/fwasaff/inelastic-rebound-md}
--- todo reproducible con \lammps{} + Python en un notebook.
\end{frame}

\end{document}
